\documentclass{article}

\usepackage[utf8]{inputenc}
\usepackage[turkish, shorthands = :!]{babel}
\usepackage[T1]{fontenc}
\usepackage{lmodern}

\usepackage{fancyhdr, lipsum, mathtools}

\renewcommand{\headrulewidth}{0 pt}
\pagestyle{fancy}
\fancyhf{}
\fancyhead[C]{\thepage}

\begin{document}

%%%%%%% SAYFA 1 %%%%%%%

\lipsum[1][1-3]
% burada asla boşluk olmamalı
\begin{align*}
\text{büyüt} \quad % amsmath paketi yüklüyken \textrm yerine \text diyebiliyoruz
	& 2 x_1 + 3 x_2 \\
\text{öyle ki} \quad
	& x_1 + x_2 \leq 2 \\
	& 4 x_1 + 6 x_2 \leq 9 \\
	& x_1, x_2 \geq 0 % buraya \\ işareti konmamalı
\end{align*}
% eğer paragraf devam ediyorsa burada da boşluk olmamalı
\lipsum[1][4-5]

\bigskip
\hrule
\bigskip

% alignat'in align'dan temel farkı, rl sütunlarının sayısının parametre olarak girilme zorunluluğu yanında, bu sütun çiftleri arasında boşluk olmaması (biz özellikle koymadıkça)

\lipsum[1][1-3]
\begin{alignat*}{1}
\text{büyüt} \quad
	& 2 x_1 + 3 x_2 \\
\text{öyle ki} \quad
	& x_1 + x_2 \leq 2 \\
	& 4 x_1 + 6 x_2 \leq 9 \\
	& x_1, x_2 \geq 0  % buraya \\ işareti konmamalı
\end{alignat*}
\lipsum[1][4-5]

\bigskip
\hrule
\bigskip

\lipsum[1][1-3]
\begin{alignat*}{2}
\text{küçült} \quad
	& \sum_{i, j} w_{ij} x_{ij} \\
\text{öyle ki} \quad
	& \sum_j x_{ij} = 1	&\quad&	\text{her $i$ için} \\
	& \sum_i x_{ij} = 1 && 		\text{her $j$ için} \\
	& x_{ij} \geq 0		&& 		\text{her $i, j$ için}
\end{alignat*}
\lipsum[1][4-5]

\newpage

%%%%%%% SAYFA 2 %%%%%%%

\lipsum[1][1-3]% bu yüzde işareti tam burada olmalı, yoksa fazladan dikey boşluk oluşuyor
\begin{subequations} \label{eq:TAP}
\begin{alignat}{2}				
\text{TAP: \quad max \quad}
	& \sum_i c_i \left(\frac{\sum_j a_{ij} b_j x_{ij}}{\sum_j b_j x_{ij}}\right)	\label{eq:TAPobj} \\
\text{s.t. \quad}
	& \sum_i x_{ij} = 1				     			&\quad& \text{for all $j$}	\label{eq:TAPhard} \\
	& \mu - d \leq \sum_j b_j x_{ij} \leq \mu + d	&&		\text{for all $i$} 	\label{eq:TAPdev} \\
	& x_{ij} \in \{0,  1\}							&& 		\text{for all $i, j$}	\label{eq:TAPvar}
\end{alignat}
\end{subequations}
\lipsum[1][4-5]

Problem~\eqref{eq:TAP}... Constraint~\eqref{eq:TAPdev}...

\end{document}